\renewcommand*{\authorOfNextLines}{Helmut Quinz}
Bei Planung und Ausführung eines Projekt gilt es einige Klippen zu umschiffen.
Häufige Fehler, welche Sie besser vermeiden  sind:
\begin{description}
	\item[Unklare Zielsetzung] Unklare oder zu vage formulierte Ziele führen zu Unsicherheit und abarbeiten an nicht benötigten Themen.
	\item[Unrealistische Zeitplanung] Eine zu kanppe Zeitplanung führt zu Stress, Verzögerungen und Qualitätsverlust.
	\item[Mangelnde Ressourcenplanung] Ein Mangel an Zeit, Geld, spezifisches Materialien oder erforderlichen Know How, wird Fortschritt im Projekt behindern
	\item[Mangelnde Kommunikation] Unzureichender Austausch unter den Team-Mitgliedern, mit dem Betreuer oder Projektpartner führt zu Missverständnissen und Fehlern.
				Legen Sie einen klaren Kommunikationsplan fest, halten Sie regelmäßige Besprechungen ab und informieren Sie alle Beteiligten über den Projektfortschritt
	\item[Perfektionismus] Der Versuch, alles perfekt zu machen, führt zu übermäßigem Zeitverbrauch und Verzögerungen im Projektverlauf.
	\item[Fehlende Flexibilität] Starre Pläne ohne Anpassungsfähigkeit führen zu Problemen bei unvorhergesehenen Ereignissen.
	
\end{description}
und besonders gefährlich:
\begin{description}
	\item[Unzureichende Selbstorganisation] führt zu einem chaotischen Arbeitsablauf und ungenügenden Fortschritt.
\end{description}